\documentclass[11pt]{article}
\title{How the Codynamic Book Machine Works}
\author{Arthur Petron}
\usepackage[margin=1in]{geometry}
\usepackage{graphicx}
\usepackage{amsmath,amssymb,mathtools}
\usepackage{tikz}
\usetikzlibrary{positioning,arrows.meta}
\usepackage{hyperref}
\graphicspath{{media/}{media/diagrams/}{images/}{artwork/}}

\begin{document}

\section*{The Repository as the Case Study}

This paper is grounded in the repository that is currently producing it. Rather than describing an abstract workflow in isolation, the manuscript treats the live project as its case study: the YAML outline that defines the work, the agent definitions that interpret it, the runtime logs that record coordination, the generated diagrams that externalize structure, and the compiled \LaTeX{} artifacts that show the system's output in a form suitable for review and publication.

That choice matters because the repository is not merely a storage location for source files. It is the operational record of the Codynamic Book Machine. The outline captures authorial intent in a machine-readable form; the section payloads translate that intent into prose; the task queues and message traces show how work is assigned and revised; and the build artifacts demonstrate whether the manuscript can actually be assembled into a coherent document. In this sense, the repository is both evidence and implementation: it contains the system and also serves as the primary evidence for how the system behaves.

Using the repository as the case study also keeps the paper concrete. Claims about coordination, specialization, and feedback are not presented as hypothetical design principles alone. They are tied to durable artifacts that can be inspected directly: the structure of the outline, the prompts and callbacks exchanged among agents, the bibliography state, the diagram requests and returned media paths, and the compiled outputs that reveal whether the manuscript remains internally consistent. The result is a working example of codynamic authoring in which the book is explained by the same repository that is used to build it.

This section therefore establishes the epistemic stance of the paper. The system is described from inside its own running repository, and the repository's artifacts are treated as the authoritative record of the architecture in action. Subsequent sections will unpack the canonical work object, the agent runtime, the editorial maintenance cycle, and the specialist handoffs that make this repository legible as a book machine rather than a loose collection of scripts and drafts.

In practical terms, the case study is not a detached example chosen after the fact; it is the live project that is generating the manuscript as it is being described. That means the paper can point to the same outline tree that schedules the chapters, the same section payloads that hold the prose, the same logs that preserve task execution, and the same compiled outputs that verify the document can be assembled. The repository is therefore the paper's laboratory, archive, and proof of concept at once.

Because the case study is the running repository itself, the paper can also distinguish between what the system claims to do and what it has actually recorded doing. The outline shows intended structure, the agent definitions show assigned responsibilities, the logs show coordination in motion, and the artifacts show the state that survived compilation. That layered evidence base is what makes the account constructive rather than speculative: the architecture is explained through the traces it leaves behind.

The remainder of the paper relies on that stance. When later sections discuss canonical data models, task queues, editorial maintenance, reference support, diagram handoffs, and artifact lifecycles, they do so against the background of this repository-level evidence. The case study is therefore not an illustrative aside; it is the ground on which the rest of the argument stands.


\end{document}
